% !TEX program = xelatex
% 问题一：多点测向误差扇区及交会定位区域示意图
% 可独立编译，也可将 tikzpicture 部分复制到论文正文。
\documentclass[tikz,border=6pt]{standalone}
\usepackage{ctex}
\usetikzlibrary{arrows.meta,angles,quotes,calc}

\definecolor{targetblue}{RGB}{54,112,178}
\definecolor{sectororange}{RGB}{230,126,34}
\definecolor{regionred}{RGB}{192,57,43}
\definecolor{stationgreen}{RGB}{39,130,92}

\begin{document}
\begin{tikzpicture}[
  x=1.55cm,
  y=1.55cm,
  >=Stealth,
  font=\small,
  boundary/.style={sectororange,dashed,line width=0.85pt},
  bearing/.style={sectororange,-{Stealth[length=2mm]},line width=0.8pt},
  station/.style={circle,fill=stationgreen,draw=white,line width=0.7pt,
                  minimum size=5.2pt,inner sep=0pt},
]
  % 比例：图中 1 个坐标单位代表 500 m。
  \def\R{3.6}
  \coordinate (O)  at (0,0);
  \coordinate (S1) at (-3.0,-1.5);
  \coordinate (S2) at ( 2.8,-1.6);
  \coordinate (S3) at (-0.5, 3.0);
  \coordinate (J)  at ( 0.3, 0.2);

  % 目标圆域。
  \fill[targetblue!4] (O) circle[radius=\R];
  \draw[targetblue,line width=1.0pt] (O) circle[radius=\R];
  \node[targetblue,anchor=south west] at (2.15,2.72)
    {目标圆域 $x^2+y^2\leq 1800^2$};

  % 三个 ±1° 扇区与目标圆域的公共交集。中心方向分别为
  % 27.26°, 144.25°, 285.95°，均穿过示意目标点 J。
  \begin{scope}
    \clip (O) circle[radius=\R];
    \clip (S1)--($(S1)+(26.26:8)$)--($(S1)+(28.26:8)$)--cycle;
    \clip (S2)--($(S2)+(143.25:8)$)--($(S2)+(145.25:8)$)--cycle;
    \clip (S3)--($(S3)+(284.95:8)$)--($(S3)+(286.95:8)$)--cycle;
    \fill[regionred!60] (-4,-4) rectangle (4,4);
  \end{scope}

  % S1 测向中心线与误差边界。
  \draw[boundary] (S1)--++(26.26:7.1);
  \draw[boundary] (S1)--++(28.26:7.1);
  \draw[bearing]  (S1)--++(27.26:4.25);

  % S2 测向中心线与误差边界。
  \draw[boundary] (S2)--++(143.25:7.1);
  \draw[boundary] (S2)--++(145.25:7.1);
  \draw[bearing]  (S2)--++(144.25:4.25);

  % S3 测向中心线与误差边界。
  \draw[boundary] (S3)--++(284.95:7.1);
  \draw[boundary] (S3)--++(286.95:7.1);
  \draw[bearing]  (S3)--++(285.95:4.25);

  % 检测点与标签。
  \node[station,label={[stationgreen]225:$S_1$}] at (S1) {};
  \node[station,label={[stationgreen]315:$S_2$}] at (S2) {};
  \node[station,label={[stationgreen]150:$S_3$}] at (S3) {};

  % 误差角局部标识（真实绘制为 1°）。
  \draw[sectororange,->,line width=0.6pt]
    ($(S1)+(27.26:0.55)$) arc[start angle=27.26,end angle=28.26,radius=0.55];
  \node[sectororange,anchor=west] at ($(S1)+(28.26:0.82)$) {$1^\circ$};
  \draw[sectororange,->,line width=0.6pt]
    ($(S1)+(27.26:0.68)$) arc[start angle=27.26,end angle=26.26,radius=0.68];
  \node[sectororange,anchor=north west] at ($(S1)+(26.26:0.92)$) {$1^\circ$};

  % 公共定位区域标签。
  \draw[regionred,-{Stealth[length=2mm]},line width=0.75pt]
    (1.35,0.78)--(0.38,0.26);
  \node[regionred,fill=white,fill opacity=0.86,text opacity=1,
        rounded corners=1pt,inner sep=2pt] at (1.62,0.92)
    {公共定位区域 $\Omega$};

  % 手工局部放大：只放大几何坐标，不放大线宽。
  \draw[regionred!75,line width=0.55pt] (J)--(-1.25,1.15);
  \begin{scope}[shift={(-2.05,1.72)},scale=9]
    \clip (0,0) circle[radius=0.115];
    \fill[white] (-0.13,-0.13) rectangle (0.13,0.13);
    \begin{scope}
      \clip (-3.3,-1.7)--($(-3.3,-1.7)+(26.26:8)$)--($(-3.3,-1.7)+(28.26:8)$)--cycle;
      \clip ( 2.5,-1.8)--($( 2.5,-1.8)+(143.25:8)$)--($( 2.5,-1.8)+(145.25:8)$)--cycle;
      \clip (-0.8, 2.8)--($(-0.8, 2.8)+(284.95:8)$)--($(-0.8, 2.8)+(286.95:8)$)--cycle;
      \fill[regionred!65] (-0.13,-0.13) rectangle (0.13,0.13);
    \end{scope}
    \draw[boundary] (-3.3,-1.7)--++(26.26:8);
    \draw[boundary] (-3.3,-1.7)--++(28.26:8);
    \draw[boundary] ( 2.5,-1.8)--++(143.25:8);
    \draw[boundary] ( 2.5,-1.8)--++(145.25:8);
    \draw[boundary] (-0.8, 2.8)--++(284.95:8);
    \draw[boundary] (-0.8, 2.8)--++(286.95:8);
  \end{scope}
  \draw[regionred,line width=0.8pt] (-2.05,1.72) circle[radius=1.035];
  \node[regionred,fill=white,inner sep=1.5pt] at (-2.05,0.53)
    {$\Omega$ 局部放大 $9\times$};

  % 坐标方向与比例尺。
  \draw[gray!75,-{Stealth[length=1.8mm]}] (-3.35,-3.12)--(-2.45,-3.12)
    node[right] {$x$（东）};
  \draw[gray!75,-{Stealth[length=1.8mm]}] (-3.35,-3.12)--(-3.35,-2.22)
    node[above] {$y$（北）};
  \draw[black,line width=1.2pt] (1.65,-3.18)--(2.65,-3.18);
  \draw[black] (1.65,-3.12)--(1.65,-3.24) (2.65,-3.12)--(2.65,-3.24);
  \node[below] at (2.15,-3.22) {$500\,\mathrm{m}$};

  % 图例。
  \begin{scope}[shift={(-0.55,-3.03)}]
    \draw[bearing] (-1.15,0)--(-0.55,0);
    \node[anchor=west] at (-0.48,0) {示向度中心线};
    \draw[boundary] (0.85,0)--(1.45,0);
    \node[anchor=west] at (1.52,0) {$\theta_i\pm1^\circ$ 边界};
  \end{scope}
\end{tikzpicture}
\end{document}
